%% thesis.tex (based on thesisPhD.tex by https://github.com/agabardo)
%% Philipp V. Rouast
%% https://github.com/prouast/uon_latex_thesis_template

% Font size and style must be reasonable for the reader to read the thesis on a screen or in hard-copy. Styles such as Arial and Times New Roman, with a font size of 10-12 are commonly used.
% The main argument of the thesis must be typed in 1.5 or double spacing. Variation in spacing may considered desirable for the presentation of tables, quotations, etc.
% The margin on each page should be not less than 4 cm on the left, 2 cm on the right, 3 cm at the top and 2 cm at the bottom for right hand pages. For left hand pages the side margins are reversed to 2cm on the left and 4 cm on the right to allow for binding.
% In cases where are hard copy submission is required, pages in the body of the thesis may be double- sided (i.e. printed on both sides) at your discretion.

% oneandhalfspace
\documentclass[a4paper,10pt,phdthesis,oneside,doublespace,pdflatex]{assets/latex/cssethesis}

% Change "phdthesis" to "confirmationreport" for confirmation report

% * Include the option "pdflatex" above if you want to use pdflatex rather than dvips
% * Include the option "oneside" if you don't want formatting for two-sided  printing.
% * include ::
%                 = singlespace}{\AtBeginDocument{\singlespacing}}
%                 = oneandhalfspace}{\AtBeginDocument{\onehalfspacing}}
%                 = doublespace}{\AtBeginDocument{\doublespacing}}
% * Include the option "nocoursecode" so that the numerical course code is suppressed after the course name.
% * Include option "thesisdraft" to get a timestamp and "Draft" message in the footer

% *** CITATION PACKAGES ***
%
\usepackage{cite}

% *** FONT/LANGUAGE PACKAGES ***
%
\usepackage[british]{babel}
\usepackage[sfdefault,light]{roboto}  %% Option 'sfdefault' only if the base font of the document is to be sans serif
\usepackage[T1]{fontenc}
\usepackage[babel=true]{microtype}
\usepackage{hyphenat}
\hyphenation{Girvan}
\usepackage{lipsum}
\usepackage{url}
\usepackage{soul} % For highlighting in draft

% *** ALGORITHM PACKAGES ***
%
\usepackage[ruled,linesnumbered]{algorithm2e} % For algorithm

% *** GRAPHICS RELATED PACKAGES ***
%
\usepackage{tikz}
\usetikzlibrary{positioning}
\usepackage[pdftex]{graphicx}
\graphicspath{{figures/}}
\DeclareGraphicsExtensions{.pdf,.png}
\usepackage{pgfplots} % For graphs
\usepgfplotslibrary{fillbetween} % For plots
\usepackage[final,poster=last]{animate} % For animations
%\onehalfspacing %%% <- Increase spacing for the proofread

% *** TYPESETTING PACKAGES ***s
%
\usepackage{setspace} %%% for additional spacing
\usepackage{xcolor}
\usepackage{framed}
\colorlet{shadecolor}{gray!20}
\usepackage{textcomp}
\makeatletter
\let\MYcaption\@makecaption
\makeatother
\usepackage[font=footnotesize]{subcaption}
\makeatletter
\let\@makecaption\MYcaption
\makeatother
\usepackage{enumitem}
\usepackage{glossaries}

% *** TABLE RELATED PACKAGES ***
%
\usepackage{array} % For bold table hline
\usepackage{booktabs}
\usepackage{tabularx}
\usepackage{pdflscape}
\usepackage{multirow}
\usepackage[flushleft]{threeparttable}
\newcolumntype{P}[1]{>{\centering\arraybackslash}p{#1}} % For centering in table
\newcolumntype{C}{>{\Centering\arraybackslash}X} % 
\usepackage{mathtools} % For matrix in table

% *** MATH RELATED PACKAGES ***
%
\usepackage{amsmath,amssymb,amsfonts}
\usepackage{float}
%

% Redefine @makechapterhead
% This is done to get rid of "Chapter X" taking up space at the start of each section.
%
\makeatletter
\def\@makechapterhead#1{%
  \vspace*{50\p@}%
  {\parindent \z@ \raggedright \normalfont
    \ifnum \c@secnumdepth >\m@ne
      \if@mainmatter
        %\huge\bfseries \@chapapp\space \thechapter
        \Huge\bfseries \thechapter.\space%
        %\par\nobreak
        %\vskip 20\p@
      \fi
    \fi
    \interlinepenalty\@M
    \Huge \bfseries #1\par\nobreak
    \vskip 40\p@
  }}
\makeatother

% Define Acronyms
\newacronym{dnn}{DNN}{deep neural network}

\makeglossaries

% *** BEGIN DOCUMENT ***
%
\begin{document}
%

% Author
\thesisauthor{Chuck Norris}
\thesisauthorlastname{Norris}
%\thesisauthorpreviousdegrees{MSc. Kung Fu.} % Optional
\thesisfield{Ass Kicking}

% Department
\thesisdepartment{\bfseries The University of Newcastle \\ Faculty of Engineering and Built Environment\\  School of Electrical Engineering and Computing}

% Date
\thesismonth{August} % Optional. Current month is used if this is not set
\thesisyear{2020} % Optional. Current year is used if this is not set

% Title
\thesistitle{Latex thesis template}

% Supervisors
\thesissupervisor{Prof Dr. Al Caholic}
\thesissupervisoremail{al.caholic@newcastle.edu.au}
\thesisassocsupervisor{Dr. Seymour Butts}
\thesisassocsupervisoremail{seymour.butts@newcastle.edu.au}

% Make front matter
\frontmatter

% The title page
% List your thesis title, your name in full, previous qualifications held in abbreviated form e.g. BSc(Hons)(Newcastle), the full name of the degree for which your thesis is submitted, and the month and year of submission of the thesis for examination.
\thesistitlepage

% Declaration pages
% After the title page, you must include a page with the following signed statement.
\thesisdeclarationonepage
\thesispublications{%
\begin{itemize}
	\item Paper 1 details
	\item Paper 2 details
	\item Paper 3 details
	\item ...
\end{itemize}
}
\thesisdeclarationthreepage

% Acknowledgments
% All candidates in receipt of a RTP Tuition Scholarship and/or RTP Fees Offset must include acknowledgement of the support received from the Australian Government Research Training Program Scholarship. All domestic students receive some support from this scholarship scheme in the form of a fees-offset scholarship.
\begin{thesisacknowledgments}
% !TEX root = thesisPhD.tex 

I here acknowledge my advisors, and everybody else that I need to acknowledge. Thank you.

\end{thesisacknowledgments}

% Table of contents
% The table of contents should provide a list of thesis chapters and major sections, with page numbers. Pages should be numbered consecutively throughout the document.
\tableofcontents

% Abstract
% The abstract should consist of approximately 300 words.
\begin{thesisabstract}
% !TEX root = thesisPhD.tex

The thesis abstract text goes here.
	
\end{thesisabstract}

% Main matter
\mainmatter

% A Doctoral thesis should not normally exceed 100,000 words, excluding appendices, tables and illustrative matter.

% Chapter 1
\glsresetall % Reset all acronyms to be written out on first use
\chapter{Introduction}
\label{ch:chapter_1}
% !TEX root = thesisPhD.tex 
LaTeX is a document preparation system for high-quality typesetting. It is most often used for medium-to-large technical or scientific documents but it can be used for almost any form of publishing.

\section{An example of text and a list}

LaTeX is based on Donald E. Knuth's TeX typesetting language or certain extensions. LaTeX was first developed in 1985 by Leslie Lamport, and is now being maintained and developed by the LaTeX3 Project. LaTeX is available for free by anonymous ftp \footnote{https://latex-project.org/ftp.html}. LaTeX contains features for:

\begin{itemize}
\item Typesetting journal articles, technical reports, books, and slide presentations.
\item Control over large documents containing sectioning, cross-references, tables and figures.
\item Typesetting of complex mathematical formulas.
\item Advanced typesetting of mathematics with AMS-LaTeX.
\item Automatic generation of bibliographies and indexes.
\item Multi-lingual typesetting.
\item Inclusion of artwork, and process or spot colour.
\item Using PostScript or Metafont fonts.
\end{itemize}

\section{An example of an image inserted}
LaTeX is a powerful typesetting system, used for producing scientific and mathematical documents of high typographic quality. Unlike WYSIWYG tools such as FrameMaker and Word, it uses plain text files that contain formatting
commands. It's big, open source, stable and used by many technical publishing companies. It's also relatively unknown in the technical writing community. This article overviews LaTeX, and directs you to sources of information. Figure \ref{figure:latex-logo} shows the beautiful  logo.

\begin{figure}[htbp]
\begin{center}
\includegraphics[width=3in]{figures/LaTeX_logo}
\caption{The LaTex logo}
\label{figure:latex-logo}
\end{center}
\end{figure}

Avoid using .JPEG, JPG, GIF, BMP or any other bitmap file. Always prefer to use .SVG or .PDF vector files when possible, the result will be much better. After inserting your figure, just continue to writing your text.

% Chapter 2
\glsresetall % Reset all acronyms to be written out on first use
\chapter{Something else}
\label{ch:chapter_2}
% !TEX root = thesisPhD.tex 
An small example using tables and algorithms.
\section{Tables}

Table \ref{table:benchmarks} is an example of table. Play around with the arraystretch parameter to get more or less space between rows. The size is set to textwidht, you can also use amounts in inches or centimeters.

\begin{table}[ht]
\renewcommand{\arraystretch}{1.1}
\centering
\resizebox{\textwidth}{!}{%
\begin{tabular}{llcccc}
\toprule
\textbf{Network}  & \textbf{Class}          & \textbf{Nodes} & \textbf{Edges} & \textbf{AVG Degree} & \textbf{Max Modularity} \\ \toprule
Zakary's Karate Club      & \small Real world network      & 34             & 78             & 4.59       & 0.4198        \\ \hline
Lusseau's Dolphins        & \small Real world network      & 62             & 159            & 5.13       & 0.5285        \\ \hline
American College Football & \small Real world network      & 115            & 613            & 10.66      & 0.6046        \\ \hline
Jazz  Musicians                    & \small Real world network      & 119            & 2742           & 27.70        & 0.4451       \\ \bottomrule
\end{tabular}
}
\caption{List of benchmark networks.}
\label{table:benchmarks}
\end{table}

\section{Algorithms}
And here, just an small example of an algorithm.

\RestyleAlgo{boxed} %Comment this line with '%' to change the aspect of the box. :-)
\begin{algorithm}[H]
\RestyleAlgo{boxruled}
\label{tourment-algorithm}
\SetAlgoLined
\DontPrintSemicolon % Some LaTeX compilers require you to use \dontprintsemicolon instead
%\KwIn{a complex network $G$.}
%\KwOut{the dendrogram with partitions of $G$.}
$Input: (Population_0, Tournment\_Size)$;\;
%$Population_0 \leftarrow Current\_Population;$\;
%$Tournment\_Size \leftarrow Tournment\_Size\_Parameter;$\;
$Population_1 = New\_Population;$\;
$Population_1 \leftarrow \{Indindividual_1,...,Individual_{Tournment\_Size}\};$\;
\For{$i =0; \leqslant Tournment\_Size$}{
	$Population_1[i] \leftarrow Population_0[random[0,Tournment\_Size]];$\;
}
$Output : Population_1.Best.$\;
\caption{Tournament selection.}
\end{algorithm}

\vspace{1cm} %This is used to add some space after the algorithm.

\paragraph{Citations:} and here, is how do you cite with a single reference \cite{harik1999compact}, or with several \cite{harik1999compact, goldberg1988genetic}. Basically, you have here all that you need to build your thesis with LaTeX. Have fun!

Using an acronym: \gls{dnn}

% ...

\appendix % all \chapter{..} commands after this will generate appendices
\glsresetall % Reset all acronyms to be written out on first use

% \chapter{Appendix chapter}
%\label{ch:appendix_a}
%\input{thesis_appendix}

\backmatter % start the thesis back matter - the numbering of chapters stops here.

% References
\pagestyle{plain}
\bibliographystyle{abbrv}
\bibliography{thesis}

% list of acronyms
\addcontentsline{toc}{chapter}{List of Acronyms}
\printglossary[type=\acronymtype,title={List of Acronyms}]

% list of figures
\listoffigures

% List of tables
\listoftables

\end{document}
